\documentclass{article}
\usepackage{amsmath}

\title{Pair Q4P9: Revelation and Contract Representation}
\author{Ada}
\date{}

\begin{document}
\maketitle

\section{The pair in two sentences}
Q studies mechanism design for agents with unknown preferences, feasible actions,
and information, and proves revelation principles that require truthful reporting
and obedience. P studies self-negotiated contracts between LLM agents and reports
a large performance difference between programmatic and natural-language contract
treatments, despite similar average use of contract-covered moves (ledger entries
6, 8, and 36).

\section{The connexion pursued}
The initial connexion claimed that natural language lets an agent reinterpret, and
thus counterfeit, a commitment. The peers rejected that claim and instead asked
whether two losslessly equivalent displays of one contract can induce different
planning behavior in an LLM agent, thereby testing the optimization and
strategy-copying abstraction used in Q's revelation construction (ledger entries
15, 32, 33, and 36).

\section{What was found}
Q's construction replaces equilibrium play in an indirect mechanism with a direct
recommendation of the induced contingent action plan. A deviation available after
the direct recommendation can be copied in the indirect mechanism, so the indirect
equilibrium blocks it; conversely, a direct incentive-compatible mechanism is
itself an indirect mechanism. Hence, under Q's assumptions, a bijective relabeling
that preserves feasible actions, information, and payoff consequences preserves
the equilibrium outcome set (ledger entries 8, 12, and 15).

P supplies suggestive behavioral evidence. On mutually dependent boards, both
agents finish in \(79\%\) of programmatic-contract games and \(46\%\) of
natural-language-contract games. Contract-covered movement is similar, about
\(2.60\) tiles per game, but natural-language games have fewer supplemental
proposals, \(1.4\) rather than \(3.0\), and agents cite the insufficient contract
in \(94\%\) rather than \(74\%\) of trade refusals (ledger entries 6 and 8).
The notes interpret this pattern as possible representation-dependent planning or
over-anchoring, not weaker enforcement.

The proposed derivation is therefore conditional. Freeze a canonical obligation
map \(c\), render the same fields and contingencies losslessly as canonical JSON
or canonical prose, and enforce both with the same exact program. At preregistered
decision nodes, compare strictly suboptimal refusals, or regret against an exact
planner. A display-dependent difference there would challenge the empirical
adequacy of unrestricted optimization and strategy copying for these agents
without refuting Q's theorem (ledger entries 23, 29, 32, and 35).

\section{Objections and their status}
The original claim fails because Q's no-counterfeiting condition concerns hard
evidence about physical execution or information capability: a certificate may
be hidden but not fabricated. In P, an external judge interprets the retained
negotiation and automatically enforces covered transfers; it makes
\(3/1155=0.26\%\) errors, all false positives. P therefore shows no
agent-controlled ex-post reinterpretation channel (ledger entries 1, 2, 33,
and 36).

The existing programmatic and natural-language treatments are not known to be
extensionally identical. They use separately generated negotiations and differ
in presentation, acceptance, retained history, and enforcement, while similar
mean coverage does not prove matched obligations (ledger entries 7, 11, 18,
and 32). Also, Q establishes equality of equilibrium outcome sets, not equality
of a particular selected outcome, so the completion-rate gap alone is not
decisive (ledger entries 21, 23, and 29). Finally, defining an
interface-dependent strategy set after observing failure would be circular; a
resource bound or decoding model must be specified or estimated independently
(ledger entries 18, 26, 32, and 35). These objections remain unresolved by the
current data and shape the proposed test.

\section{Outside sources}
No source outside Q and P was used to establish the connexion. The ledger and
the two peer notes only interpret results and arguments contained in Q and P.

\section{What remains open}
It remains unknown whether a representation effect survives matched semantics
and common enforcement. It is also open whether any surviving effect is
specific to incomplete contracts, whether it reflects context or syntax, and
whether an independently specified bounded decoder-policy model predicts it.
The material is thin on these points because no matched intervention has yet
been run (ledger entries 12, 14, 32, and 35).

\section{Smallest next step}
Compile one accepted obligation map into losslessly equivalent canonical JSON
and canonical prose, show each rendering under identical programmatic
enforcement on the same board and model, and test preregistered nodes where a
supplemental trade is strictly payoff-improving under every relevant
continuation. A representation-dependent difference in refusal or exact-planner
regret would establish the proposed behavioral gap; no difference would locate
P's original contrast in unmatched content, enforcement, or sampling rather
than representation (ledger entries 15, 23, 32, and 35).

\end{document}
